% ============================================================
% Template LaTeX ABNT NBR 14724/2024 — IFB
% Compatível com pdflatex + abntex2cite
% ============================================================

\documentclass[
  12pt,
  a4paper,
  oneside,
  brazil,
  chapter=TITLE
]{article}

% ── Codificação e Idioma ──────────────────────────────────────────────────────
\usepackage[T1]{fontenc}
\usepackage[utf8]{inputenc}
\usepackage[brazil]{babel}

% ── Fonte Times New Roman ─────────────────────────────────────────────────────
\usepackage{times}

% ── Matemática ────────────────────────────────────────────────────────────────
\usepackage{amsmath}

% ── Geometria da página — margens ABNT ────────────────────────────────────────
\usepackage[
  a4paper,
  left=3cm,
  top=3cm,
  right=2cm,
  bottom=2cm,
  headheight=14pt,
  headsep=0.5cm
]{geometry}

% ── Espaçamento 1,5 ───────────────────────────────────────────────────────────
\usepackage{setspace}
\onehalfspacing

% ── Recuo de parágrafo: 1,25 cm ───────────────────────────────────────────────
\usepackage{indentfirst}
\setlength{\parindent}{1.25cm}
\setlength{\parskip}{0pt}

% ── Formatação de seções — NBR 6024 ───────────────────────────────────────────
\usepackage{titlesec}

\titleformat{\section}
  {\normalfont\normalsize\bfseries\MakeUppercase}
  {\thesection\quad}{0em}{}

\titleformat{\subsection}
  {\normalfont\normalsize\bfseries}
  {\thesubsection\quad}{0em}{}

\titleformat{\subsubsection}
  {\normalfont\normalsize\bfseries\itshape}
  {\thesubsubsection\quad}{0em}{}

\titleformat{\paragraph}[runin]
  {\normalfont\normalsize\itshape}
  {\theparagraph\quad}{0em}{}[\hspace{0.5em}]

\titlespacing*{\section}      {0pt}{18pt}{12pt}
\titlespacing*{\subsection}   {0pt}{18pt}{6pt}
\titlespacing*{\subsubsection}{0pt}{12pt}{6pt}

% Seções primárias iniciam em nova página
\newcommand{\secbreak}{\clearpage}
\let\oldsection\section
\renewcommand{\section}{\secbreak\oldsection}

% ── Paginação — canto superior direito, 10pt ──────────────────────────────────
\usepackage{fancyhdr}
\pagestyle{fancy}
\fancyhf{}
\fancyhead[R]{\fontsize{10}{12}\selectfont\thepage}
\renewcommand{\headrulewidth}{0pt}
\renewcommand{\footrulewidth}{0pt}
\fancypagestyle{plain}{
  \fancyhf{}
  \fancyhead[R]{\fontsize{10}{12}\selectfont\thepage}
  \renewcommand{\headrulewidth}{0pt}
}

% ── Tabelas e Quadros ─────────────────────────────────────────────────────────
\usepackage{longtable}
\usepackage{booktabs}
\usepackage{array}
\usepackage{multirow}
\usepackage{tabularx}
\renewcommand{\arraystretch}{1.3}

% ── Legendas ──────────────────────────────────────────────────────────────────
\usepackage{caption}
\DeclareCaptionFormat{abnt}{\fontsize{10}{12}\selectfont#1#2#3}
\DeclareCaptionLabelSeparator{emdash}{\,---\,}
\captionsetup{
  format=abnt,
  justification=centering,
  singlelinecheck=false,
  labelsep=emdash,
  position=above,
  skip=3pt
}
\newcommand{\fonte}[1]{%
  \begin{flushleft}\fontsize{10}{12}\selectfont Fonte: #1\end{flushleft}
}

% ── Listas ───────────────────────────────────────────────────────────────────
\usepackage{enumitem}
\setlist{nosep, leftmargin=*}

% ── Sumário e Lista de Quadros ────────────────────────────────────────────────
\usepackage{tocloft}
\renewcommand{\contentsname}{SUMÁRIO}
\newlistof{quadros}{loq}{LISTA DE QUADROS}
\setlength{\cftquadrosindent}{0pt}
\setlength{\cftquadrosnumwidth}{0pt}
\makeatletter
\renewcommand{\listofquadros}{%
  \section*{LISTA DE QUADROS}%
  \@starttoc{loq}%
}
\makeatother

\newcommand{\quadrotitle}[2]{%
  \vspace{1em}%
  \phantomsection\label{quadro#1}%
  \addcontentsline{loq}{quadros}{Quadro #1 --- #2}%
  \noindent\textbf{Quadro #1 --- #2}\par\vspace{0.3em}%
}
\newcommand{\quadrofonte}[1]{\par\noindent\textit{Fonte: #1}\vspace{1em}\par}

% ── Hyperlinks ────────────────────────────────────────────────────────────────
\PassOptionsToPackage{hyphens}{url}
\usepackage[hidelinks, colorlinks=false, breaklinks=true]{hyperref}

% ── Citações longas ───────────────────────────────────────────────────────────
\usepackage{quoting}
\quotingsetup{font=small, indentfirst=false, vskip=\baselineskip}

% ── Figuras ───────────────────────────────────────────────────────────────────
\usepackage{graphicx}
\usepackage{float}

% ── URL ──────────────────────────────────────────────────────────────────────
\usepackage{url}

\date{}

% ═════════════════════════════════════════════════════════════════════════════
\begin{document}

% ─────────────────────────────────────────────────────────────────────────────
% CAPA — NBR 14724/2024
% ─────────────────────────────────────────────────────────────────────────────
\begin{titlepage}
  \thispagestyle{empty}
  \begin{center}

    \includegraphics[width=0.25\textwidth]{../../logo_ifb_despro.png}\\[1.8em]

    {\normalsize\bfseries
      Instituto Federal de Educação, Ciência e Tecnologia de Brasília --- IFB\\[0.3em]
      \textit{Campus} Samambaia\\[0.3em]
      Curso Superior de Tecnologia em Design de Produto
    }

    \vfill

    {\large\bfseries
      PROPRIEDADES ORGANOLÉPTICAS DE OBJETOS SELECIONADOS
    }\\[1.2em]

    {\normalsize
      Análise sensorial aplicada a materiais, processos e experiência de uso
    }

    \vfill

    {\normalsize Brasília/DF\\2026}

  \end{center}
\end{titlepage}
\setcounter{page}{2}

% ─────────────────────────────────────────────────────────────────────────────
% FOLHA DE ROSTO — NBR 14724/2024
% ─────────────────────────────────────────────────────────────────────────────
\begin{titlepage}
  \thispagestyle{empty}
  \begin{center}

    {\normalsize\bfseries VICTOR HUGO DA SILVA COSTA}\\[0.3em]
    {\normalsize Matrícula: 2026103970002}

    \vfill

    {\large\bfseries
      PROPRIEDADES ORGANOLÉPTICAS DE OBJETOS SELECIONADOS
    }

    \vfill

    \begin{flushright}
      \begin{minipage}{0.5\textwidth}
        \normalsize
        Relatório apresentado ao Curso Superior de
        Tecnologia em Design de Produto do Instituto
        Federal de Educação, Ciência e Tecnologia de
        Brasília -- Campus Samambaia, como requisito
        parcial para a aprovação na disciplina de
        Tecnologia e Propriedade dos Materiais.

        \vspace{1em}
        Professora: Profª. Keila Sanches
      \end{minipage}
    \end{flushright}

    \vfill

    {\normalsize Brasília/DF -- 2026}

  \end{center}
\end{titlepage}

% ─────────────────────────────────────────────────────────────────────────────
% RESUMO — NBR 6028/2003
% ─────────────────────────────────────────────────────────────────────────────
\clearpage
\thispagestyle{empty}
\begin{center}
  {\normalsize\bfseries RESUMO}
\end{center}

\vspace{1em}
\noindent
Este trabalho analisa as propriedades organolépticas de dez objetos de design de produto,
organizados em pares por material predominante: alumínio anodizado, vidro tratado,
silicone elastomérico (LSR), ABS e cerâmica técnica --- sob a perspectiva da experiência
sensorial do usuário. A partir de grandezas físicas mensuráveis como efusividade térmica,
coeficiente de atrito e módulo elástico, investiga-se como os processos de fabricação e
acabamento configuram a superfície e determinam a percepção tátil, visual, acústica e,
em alguns casos, olfativa dos objetos. A análise evidencia que a hierarquia sensorial de um
produto emerge da interação entre materiais distintos e que a especificação das propriedades
organolépticas deve integrar o processo projetual desde as fases iniciais, não como etapa
de acabamento.

\vspace{1em}
\noindent\textbf{Palavras-chave:} propriedades organolépticas; materiais; design de produto;
efusividade térmica; experiência sensorial.

% ─────────────────────────────────────────────────────────────────────────────
% LISTA DE QUADROS
% ─────────────────────────────────────────────────────────────────────────────
\clearpage
\thispagestyle{empty}
\listofquadros

% ─────────────────────────────────────────────────────────────────────────────
% SUMÁRIO
% ─────────────────────────────────────────────────────────────────────────────
\clearpage
\thispagestyle{empty}
\tableofcontents

% ═════════════════════════════════════════════════════════════════════════════
% CONTEÚDO
% ═════════════════════════════════════════════════════════════════════════════

\section{Introdução}

A seleção de materiais em design de produto é tradicionalmente orientada por critérios de
desempenho mecânico, custo e processabilidade. Resistência à tração, módulo elástico,
temperatura de deflexão: esses números dominam as fichas técnicas e os fluxos de decisão
projetual. O problema é que o usuário não consulta fichas técnicas: ele pega o produto,
passa o dedo, sente o peso, ouve o clique. A experiência que define se um objeto parece
bom ou parece barato é, em grande medida, organoléptica.

Propriedades organolépticas são aquelas percebidas pelos órgãos dos sentidos: o que se vê,
toca, ouve e, eventualmente, cheira ou prova. Aplicadas a materiais de engenharia, essas
propriedades emergem de grandezas físicas mensuráveis e de como os processos de fabricação
configuram a superfície do objeto. Microtextura, tratamentos térmicos e químicos,
revestimentos e porosidade controlada são os recursos que o designer ajusta para modular
a experiência sensorial.

Este trabalho analisa dez objetos de design de produto, organizados em pares por material
predominante: dois em alumínio anodizado (corpo de smartphone e maçaneta de porta), dois
em vidro tratado (copo de borossilicato e tela de smartphone), dois em silicone elastomérico
LSR (bico de mamadeira e squeeze de ciclismo), dois em ABS (mouse e controle remoto) e dois
em cerâmica técnica (caneca de porcelana e caixa de relógio em zircônia). Cada par permite
comparar como o mesmo material se comporta em contextos de uso radicalmente diferentes,
o que revela mais sobre a lógica organoléptica do design do que a análise isolada de cada
material conseguiria.


\section{Mecanismos de percepção organoléptica}

A percepção sensorial de materiais não é arbitrária. Três mecanismos físicos subjacentes
explicam a maior parte do que sentimos ao tocar um objeto, e convém conhecê-los antes de
examinar os casos.

A troca térmica na pele talvez seja o mecanismo mais contraintuitivo. A sensação de
``frio'' ou ``quente'' ao primeiro toque não depende da temperatura do objeto, mas da
velocidade com que ele extrai calor da mão, governada pela efusividade térmica:
$e = \sqrt{k \rho c_p}$, onde $k$ é condutividade, $\rho$ é densidade e $c_p$ é calor
específico. Materiais com alta efusividade, como o alumínio, são percebidos como frios
e ``premium''. Materiais com baixa efusividade, como o silicone, parecem quentes ou
neutros (ASHBY; JOHNSON, 2010).

Atrito e área real de contato explicam a sensação de ``sedoso'', ``escorregadio'' ou
``grudento''. A interação entre pele e topografia superficial varia com microtextura,
umidade e coatings. Pequenas mudanças de superfície alteram o coeficiente de atrito (CoF)
de forma mensurável e, com ele, as avaliações subjetivas de conforto.

Rigidez, massa e amortecimento, por sua vez, determina o vocabulário acústico do objeto.
Módulos elásticos altos e densidades elevadas produzem transientes curtos e agudos;
polímeros e elastômeros amortecem e ``matam'' o som. É por isso que vidro e cerâmica soam
tão diferentes de ABS, mesmo quando a geometria é idêntica.

O bulk do material define a ordem de grandeza dessas propriedades, mas é o processo de
superfície que decide a experiência final. A mesma liga de alumínio pode parecer sedosa
ou áspera conforme o acabamento.


\section{Análise dos objetos}

\subsection{Corpo de smartphone (alumínio anodizado)}

O unibody de alumínio anodizado é a referência contemporânea do ``frio premium'': toque
imediato que comunica rigidez, precisão e valor. O escovamento direcional evidencia
usinagem mecânica e orienta o percurso do olhar sobre a forma; a anodização permite
coloração sem perder a leitura metálica. Acusticamente, a resposta é seca, com transientes
curtos que reforçam a percepção de solidez quando o conjunto está bem rigidizado.

Para a liga 6061, o NIST registra $k \approx 154\ \mathrm{W/(m \cdot K)}$ e
$c_p \approx 943\ \mathrm{J/(kg \cdot K)}$, resultando em
$e \approx 2{,}0 \times 10^{4}\ \mathrm{W \cdot s^{0{,}5} \cdot m^{-2} \cdot K^{-1}}$,
valor muito superior ao dos demais materiais analisados e que explica a intensidade da
sensação fria ao primeiro toque. O módulo elástico de 69--70 GPa e a espessura controlada
sustentam a percepção de precisão e ausência de flexão.

A anodização Tipo III (hard anodize, MIL-A-8625) adiciona camada de óxido com microdurezas
de centenas de HV. Selagem produz toque ``fechado e limpo''; sem selagem, há risco de
\textit{chalky feel}. Coatings oleofóbicos (AF) completam o sistema ao reduzir marcas
digitais em zonas de alto toque.

\subsection{Maçaneta de porta (alumínio anodizado)}

Na maçaneta, as propriedades organolépticas do alumínio são testadas numa condição
diferente: toque repetido, curto e funcional, sem o contexto de ``objeto de valor'' do
smartphone. O que muda não é o material, é o que o usuário traz para o contato.

A efusividade alta
(${\sim}2 \times 10^{4}\ \mathrm{W \cdot s^{0{,}5} \cdot m^{-2} \cdot K^{-1}}$)
produz sensação fria que pode oscilar entre ``qualidade'' e ``desconforto'' dependendo da
temperatura ambiente e da duração do contato. Jateamento isotrópico cria grip suave e
uniforme, adequado ao giro da mão. A geometria tem papel acústico: espessuras bem
calibradas produzem som seco ao impacto, reforçando a percepção de solidez.

O contraste entre o alumínio frio de uma maçaneta e a madeira quente de uma porta é um
exemplo direto de como a hierarquia sensorial se constrói por diferença de efusividade,
não por escolha decorativa: trata-se de um fenômeno físico cotidiano cujo fundamento
raramente é explicitado no processo projetual.

\subsection{Copo de borossilicato}

O vidro borossilicato combina transparência óptica profunda com estabilidade térmica
superior ao vidro soda-lime comum ($k \approx 1{,}2\ \mathrm{W/(m \cdot K)}$,
$c_p \approx 830\ \mathrm{J/(kg \cdot K)}$,
$e \approx 1{,}6 \times 10^{3}\ \mathrm{W \cdot s^{0{,}5} \cdot m^{-2} \cdot K^{-1}}$).
O módulo elástico elevado (~64 GPa) confere ao material um transiente acústico agudo e
característico ao impacto, percebido como sinal de rigidez e densidade.

O acabamento de borda é talvez o detalhe mais revelador desse objeto: borda polida transmite
precisão e segurança ao lábio; borda cortante ou irregular destrói a experiência
independentemente de qualquer outra qualidade. Transparência total permite leitura visual
do conteúdo, o que é tanto estética quanto função no contexto alimentar. O vidro não afeta
sabor nem odor, ao contrário de plásticos e alguns metais.

\subsection{Tela de smartphone (vidro quimicamente reforçado)}

A tela de smartphone é provavelmente a superfície mais tocada de qualquer produto
contemporâneo. Para vidro aluminosilicato reforçado, a compressão superficial pode atingir
${\sim}950\ \mathrm{MPa}$ com profundidade de camada de dezenas de micrômetros
(CORNING, [s.d.]), elevando resistência a riscos e dureza por indentação.

O CoF na interação dedo-vidro varia entre 0,7 e 1,2 conforme textura e umidade da pele.
A especificação de coatings oleofóbicos, que reduzem marcas e modulam o glide, integra
portanto a engenharia de superfície à experiência gestual do usuário.

Juhani Pallasmaa, em ``Os Olhos da Pele'' (2011), argumenta que a arquitetura contemporânea
privilegia a visão em detrimento dos demais sentidos, produzindo espaços que se deixam
fotografar mas não convidam o corpo a ficar. A tela de vidro literaliza essa crítica no
objeto de bolso: toda a riqueza sensorial do gesto foi reduzida a uma superfície plana e
lisa por onde os olhos passam enquanto os dedos deslizam sobre nada. O som de toque também
comunica qualidade; vidro mal montado ressoa diferente de vidro bem assentado em estrutura
rígida.

\subsection{Bico de mamadeira (silicone LSR)}

O bico de mamadeira é o objeto desta análise em que a dimensão organoléptica mais crítica
é aquela que não se vê: cheiro e sabor. Graus para contato oral exigem conformidade
USP Class VI e regulamentações de contato alimentar, com controle rigoroso de voláteis,
extrativos e catalisadores residuais. Silicone de grau inadequado pode ter odor residual
de catalisadores, e isso encerra a experiência imediatamente.

A efusividade baixa
($e \approx 5 \times 10^{2}\ \mathrm{W \cdot s^{0{,}5} \cdot m^{-2} \cdot K^{-1}}$)
produz sensação neutra, ideal para contato prolongado. A dureza Shore A define snap-back e
resistência ao colapso durante a sucção. A textura do molde é o único acabamento disponível,
pois silicone LSR sai da injeção pronto.

\subsection{Squeeze de ciclismo (silicone LSR)}

No squeeze, o silicone opera sob esforço mecânico repetido e condições variáveis de
temperatura e suor. A maciez (Shore~A ${\sim}$40--50) permite compressão com uma mão;
o retorno elástico devolve o volume original. Grip em superfície molhada é função direta
da microtextura do molde: texturas em diamante ou micro-canais aumentam CoF mesmo com
luvas. O mate absorve luz, coerente com o contexto esportivo, e o material não transmite
vibração para a mão durante o uso.

\subsection{Mouse (ABS)}

O mouse é um dos testes mais exigentes para ABS: superfície de médio a alto toque
contínuo, onde a percepção de solidez depende quase inteiramente da geometria. O módulo
de tração (${\sim}2.400\ \mathrm{MPa}$, SABIC CYCOLAC MG94F) é cerca de trinta vezes
menor que o do alumínio, o que significa que nervuras, espessura local e encaixes
determinam se o objeto parece sólido ou cede ao aperto.

A efusividade entre $5 \times 10^{2}$ e $8 \times 10^{2}\ \mathrm{W \cdot s^{0{,}5}
\cdot m^{-2} \cdot K^{-1}}$ sustenta toque neutro. Texturas difusas no molde reduzem
marcas de dedo sem coating dedicado. O som do clique, um dos fatores mais determinantes
de qualidade percebida em periféricos, é função do conjunto switch + carcaça + rigidez
do conjunto; o ABS amortece menos que metal, o que pode ser positivo ou negativo conforme
o perfil de som desejado.

\subsection{Controle remoto (ABS)}

No controle remoto, o ABS opera num contexto de atenção mínima: o objeto é manuseado
brevemente e com baixo engajamento visual. A qualidade percebida depende de três fatores
organolépticos principais: o som dos botões, a textura da carcaça e o peso --- objetos
com massa insuficiente tendem a ser avaliados negativamente em termos de durabilidade
percebida.

Pintura mate sobre ABS reduz marcas de dedo. UV hardcoat eleva dureza, mas o desgaste
localizado ao longo do tempo expõe o plástico cru abaixo e é percebido imediatamente como
queda de qualidade. A geometria dos botões, especialmente profundidade de curso e retorno
tátil, é tão determinante quanto o material da carcaça.

\subsection{Caneca de porcelana (cerâmica esmaltada)}

A caneca de porcelana é o objeto desta análise em que as propriedades organolépticas
cerâmicas se apresentam de forma mais completa, mobilizando simultaneamente cinco
modalidades sensoriais: visual (brilho profundo do esmalte), tátil (superfície lisa e fria),
acústico (transiente metálico ao contato com utensílios), olfativo (inércia química que
não interfere no aroma) e gustativo (ausência de migração de compostos para o líquido).

A efusividade da cerâmica esmaltada situa-se entre $2{,}4 \times 10^{3}$ e
$3{,}1 \times 10^{3}\ \mathrm{W \cdot s^{0{,}5} \cdot m^{-2} \cdot K^{-1}}$, mais fria
que polímeros com sensação distinta do vidro. O esmalte adiciona camada vítrea com cor,
brilho e resistência química; dureza Mohs mínima ${\sim}5$ garante resistência a riscos
cotidianos. A estabilidade cromática do esmalte é superior à de tintas orgânicas: a cor
de uma caneca de porcelana não desvanece.

\subsection{Caixa de relógio (cerâmica técnica --- zircônia)}

A caixa em zircônia técnica é o extremo superior da hierarquia organoléptica entre os dez
objetos. Módulo elástico ${\sim}200\ \mathrm{GPa}$, dureza ${\sim}1.300\ \mathrm{HV}$,
rugosidade superficial na escala nanométrica após polimento diamantado. O resultado é um
toque que a indústria relojoeira descreve como ``joia'': liso de uma forma que nenhum
metal ou vidro consegue reproduzir exatamente.

A efusividade entre $2{,}4 \times 10^{3}$ e
$3{,}1 \times 10^{3}\ \mathrm{W \cdot s^{0{,}5} \cdot m^{-2} \cdot K^{-1}}$ produz
sensação fria limpa, perceptivelmente diferente do alumínio. O ``clack'' agudo ao toque
comunica dureza excepcional. A estabilidade cromática intrínseca elimina o desbotamento
que anodização metálica pode sofrer com uso e suor. A fragilidade a impacto localizado é
a principal limitação: cantos vivos e tensões concentradas exigem atenção no projeto.


\section{Síntese comparativa}

O Quadro~1 sintetiza os dez objetos com suas assinaturas organolépticas, materiais
primários, efusividade de referência e processo de acabamento mais determinante.

\quadrotitle{1}{Assinaturas organolépticas dos dez objetos analisados}

\begin{longtable}[]{@{}
  >{\raggedright\arraybackslash}p{\dimexpr(\linewidth - 8\tabcolsep)*20/100\relax}
  >{\raggedright\arraybackslash}p{\dimexpr(\linewidth - 8\tabcolsep)*16/100\relax}
  >{\raggedright\arraybackslash}p{\dimexpr(\linewidth - 8\tabcolsep)*20/100\relax}
  >{\raggedright\arraybackslash}p{\dimexpr(\linewidth - 8\tabcolsep)*20/100\relax}
  >{\raggedright\arraybackslash}p{\dimexpr(\linewidth - 8\tabcolsep)*24/100\relax}@{}}
\toprule\noalign{}
\begin{minipage}[b]{\linewidth}\raggedright
\textbf{Objeto}
\end{minipage} &
\begin{minipage}[b]{\linewidth}\raggedright
\textbf{Material}
\end{minipage} &
\begin{minipage}[b]{\linewidth}\raggedright
\textbf{Assinatura sensorial}
\end{minipage} &
\begin{minipage}[b]{\linewidth}\raggedright
\textbf{$e$} (\small{W$\cdot$s$^{0{,}5}$\,m$^{-2}$\,K$^{-1}$})
\end{minipage} &
\begin{minipage}[b]{\linewidth}\raggedright
\textbf{Processo determinante}
\end{minipage} \\
\midrule\noalign{}
\endhead
\bottomrule\noalign{}
\endlastfoot
Corpo de smartphone & Al anodizado & Frio, rígido, técnico &
  ${\sim}2 \times 10^{4}$ &
  CNC + escovado + anodização + AF \\
Maçaneta de porta & Al anodizado & Frio funcional, grip suave &
  ${\sim}2 \times 10^{4}$ &
  Jateado + anodização \\
Copo de borossilicato & Vidro & Frio, transparente, acústico &
  ${\sim}1{,}6 \times 10^{3}$ &
  Polimento de borda \\
Tela de smartphone & Vidro reforçado & Liso, glide variável &
  ${\sim}1{,}4 \times 10^{3}$ &
  Ion exchange + coating AF \\
Bico de mamadeira & Silicone LSR & Neutro, macio, inerte &
  ${\sim}5 \times 10^{2}$ &
  LIM/LSR grau alimentar \\
Squeeze de ciclismo & Silicone LSR & Macio, grip, elástico &
  ${\sim}5 \times 10^{2}$ &
  LIM/LSR + textura de molde \\
Mouse & ABS & Neutro, geometria-dependente &
  ${\sim}5\text{--}8 \times 10^{2}$ &
  Injeção + textura + encaixes \\
Controle remoto & ABS & Neutro, leve, clique &
  ${\sim}5\text{--}8 \times 10^{2}$ &
  Injeção + pintura mate \\
Caneca de porcelana & Cerâmica esmaltada & Frio limpo, brilho profundo &
  ${\sim}2{,}4\text{--}3{,}1 \times 10^{3}$ &
  Sinterização + esmalte \\
Caixa de relógio & Cerâmica (zircônia) & ``Joia'', dureza excepcional &
  ${\sim}2{,}4\text{--}3{,}1 \times 10^{3}$ &
  CIM + polimento diamantado \\
\end{longtable}

\quadrofonte{elaborado pelo autor.}

A hierarquia de efusividade atravessa os dez objetos na mesma ordem:

\[
\mbox{Alumínio} \gg \mbox{Cerâmica} > \mbox{Vidro} \gg \mbox{ABS} \approx \mbox{Silicone}
\]

Produtos de alto valor percebido constroem sua experiência por contrastes deliberados entre
materiais. A frieza do alumínio ganha significado quando justaposta à maciez do silicone;
a precisão óptica do vidro se destaca contra a opacidade mate do ABS. Projetar esses
contrastes com intenção é o que diferencia a seleção de materiais como exercício técnico
da seleção de materiais como ato projetual (ASHBY; JOHNSON, 2010).


\section{Considerações finais}

A análise dos dez objetos revela um princípio consistente: em todos eles, a superfície é
o principal ``botão'' de organolepsia. O bulk define a ordem de grandeza das propriedades
sensoriais, mas é o pacote de processo que decide se a experiência será percebida como
premium ou genérica.

As propriedades organolépticas não podem ser tratadas como acabamento tardio. Elas precisam
integrar a especificação de materiais desde as fases iniciais do projeto com métricas
mensuráveis: efusividade, CoF, rugosidade, dureza superficial. A escolha de um material é
também a escolha de uma experiência sensorial, e o designer que compreende os mecanismos
por trás dessa experiência tem mais controle sobre o resultado do que aquele que se apoia
apenas na intuição (LÖBACH, 2001).


% ─────────────────────────────────────────────────────────────────────────────
% REFERÊNCIAS
% ─────────────────────────────────────────────────────────────────────────────
\section{Referências}\label{referencias}

\noindent
ASHBY, M. F.; JOHNSON, K. \textbf{Materials and design}: the art and science of material
selection in product design. 2. ed. Amsterdam: Elsevier/Butterworth-Heinemann, 2010.

\bigskip\noindent
ASTM INTERNATIONAL. \textbf{ASTM C1048}: Standard specification for heat-strengthened and
fully tempered flat glass. West Conshohocken: ASTM, 2018.

\bigskip\noindent
CALEGARI, E. P.; OLIVEIRA, B. F. de. Um estudo focado na relação entre design e materiais.
\textbf{Projética}, Londrina, v. 4, n. 1, p. 49--64, jan./jun. 2013.

\bigskip\noindent
CORNING INCORPORATED. \textbf{Gorilla Glass 3 product datasheet}. Corning: Corning Inc.,
[s.d.].

\bigskip\noindent
LÖBACH, Bernd. \textbf{Design industrial}: bases para a configuração dos produtos
industriais. São Paulo: Edgard Blücher, 2001.

\bigskip\noindent
MILITARY SPECIFICATION. \textbf{MIL-A-8625F}: anodic coatings for aluminum and aluminum
alloys. Washington: US Department of Defense, 1993.

\bigskip\noindent
NATIONAL INSTITUTE OF STANDARDS AND TECHNOLOGY (NIST). \textbf{Thermophysical properties
of aluminum alloy 6061}. Gaithersburg: NIST, [s.d.].

\bigskip\noindent
PALLASMAA, Juhani. \textbf{Os olhos da pele}: a arquitetura dos sentidos. Porto Alegre:
Bookman, 2011.

\bigskip\noindent
SABIC. \textbf{CYCOLAC MG94F ABS resin datasheet}. Pittsfield: SABIC, [s.d.].

\end{document}
